% Pomerenke's Conjecture
% Schwarzian Derivatives of Hyp Convex functions
% Second Revision for PLMS - Summer 2005

\documentclass{lms}
\usepackage{amsmath,amssymb,amsfonts,graphics,epsfig,rotating,cite,rotating}

\newcommand{\C}{\mathbb C}
\newcommand{\D}{\mathbb D}
\newcommand{\bdryD}{\partial \mathbb D}
\newcommand{\R}{\mathbb R}
\newcommand{\Ub}{\mathbb U} 
\newcommand{\Lb}{\mathbb L}
\newcommand{\HH}{\mathbb H}
\newcommand{\Z}{\mathbb Z}
\newcommand{\K}{\mathcal K}
\newcommand{\Kh}{\widetilde{\mathcal K}}
\newcommand{\Dh}{\widetilde{D}}
\newcommand{\Ph}{\widetilde{P}}
\newcommand{\Sh}{\widetilde{S}}
\newcommand{\Pt}{\widetilde{P}}
\newcommand{\Tt}{\widetilde{T}}
\newcommand{\tgamma}{\widetilde\gamma}
\newcommand{\ft}{\widetilde f}
\newcommand{\LL}{\mathcal L}
\newcommand{\PP}{\mathcal P}
\newcommand{\T}{\mathcal T}
\newcommand{\CC}{\mathcal C}
\newcommand{\hh}{\widehat h}
\newcommand{\tp}{{\tau^\prime}}
\newcommand{\tn}{{\tau_n}}
\newcommand{\ti}{{\tau_\infty}}
\newcommand{\Kp}{\mathcal K^\prime}
\newcommand{\Kpt}{\mathcal \widetilde{K^\prime}}
\newcommand{\Qt}{\widetilde Q}
\newcommand{\ap}{\alpha^\prime}
\newcommand{\tpn}{{\tau_n^\prime}}
\newcommand{\tphi}{\widetilde \varphi}
\newcommand{\abs}[1]{\lvert#1\rvert}
\newcommand{\norm}[1]{\lvert\lvert#1\rvert\rvert}
\newcommand{\Snf}{\norm{S_f}_\D}
\newcommand{\PSL}{\text{PSL}}
\newcommand{\z}{\zeta}
\newcommand{\e}{\epsilon}
\newcommand{\iG}{\int_{\gamma}}
\newcommand{\G}{{\scriptscriptstyle{\Gamma}}}



%\theoremstyle{plain}
\newtheorem{theorem}{Theorem}[section]
\newtheorem{lemma}[theorem]{Lemma}
\newtheorem{proposition}[theorem]{Proposition}
\newtheorem{corollary}[theorem]{Corollary}
\newtheorem{julia}{Julia Variational Formula}
\newtheorem{thm}{Theorem 1.1}
\newtheorem{step-down}{Step Down Lemma}

%\theoremstyle{definition} 
\newtheorem{definition}[theorem]{Definition}
\newtheorem{notation}[theorem]{Notation}

%\theoremstyle{remark} 
\newtheorem{remark}[theorem]{Remark}
\newtheorem{example}[theorem]{Example}
\newtheorem{note}[theorem]{Note}

\DeclareMathOperator{\m}{mod}
\DeclareMathOperator{\im}{Im}
\DeclareMathOperator{\re}{Re}
\DeclareMathOperator{\Arg}{Arg}
\DeclareMathOperator{\Log}{Log}
\DeclareMathOperator{\interior}{int}


\numberwithin{equation}{section}


\begin{document}



\title[Deformation of a Disc Under a Hyperbolically Convex Map]{The Sharp Bound
  for the Deformation of a Disc under a Hyperbolically  Convex Map}
%[Schwarzian Derivatives of Hyperbolically Convex Functions]{A Sharp 
%Bound on the Schwarzian Derivatives of Hyperbolically
%Convex Functions}

\date{\today}

\author{Roger W. Barnard, Leah Cole, Kent Pearce, and G. Brock Williams}



%\keywords{Hyperbolically Convex, Schwarzian Derivative, Julia Variation}

%\subjclass{30C70}

%Insert `2000 Mathematics Subject Classification' numbers here!
\classno{30C70}

\maketitle

\begin{abstract}
We complete the determination of how far convex maps can deform discs in 
each of the three classical geometries.
The euclidean case was settled by Nehari in 1976, and the spherical case
by Mej\'ia and Pommerenke in 2000.
We find the sharp bound on the Schwarzian derivative of a hyperbolically convex 
function and thus complete the hyperbolic case.
 This problem was first posed by Ma and Minda in a series of papers in the 1980's.
Mej\'ia and Pommerenke then produced partial results and a conjecture
as to the extremal function in 2000.  Their function maps onto a domain bounded by
two proper geodesic sides, a ``hyperbolic strip.''
Applying a generalization of the Julia variation and a critical Step Down Lemma,
 we show that 
there is an  extremal function mapping  onto a domain with 
at most two geodesic sides.  We then verify using special function theory
that among the remaining candidates, Mej\'ia and Pommerenke's  two-sided 
domain is in fact extremal.  This correlates nicely with the euclidean
and spherically convex cases in which the extremal is known to be
a map onto a two-sided ``strip.'' 
\end{abstract}




%******************************************************************

\section{Introduction}\label{S:intro}

Hyperbolic convexity is a natural generalization of euclidean
convexity; a region $G$ in the Poincar\'e model $\D$ of the hyperbolic
plane  is  \textbf{hyperbolically convex} if for any
two points in $G$, the hyperbolic geodesic segment between them lies entirely 
in $G$.  Such regions arise naturally  in Teichm\"uller theory, for example, since
the fundamental domains of Fuchsian groups are hyperbolically convex~\cite{aB83,
oL87}.

A conformal map $f:\D\to\D$ is hyperbolically convex if
its range is hyperbolically convex.
Hyperbolically convex functions have been extensively studied  by Ma and 
Minda~\cite{MM94,MM99} and Mej\'ia and Pommerenke~\cite{PM00,MP00,MP01,MP02,MP00b},
as well as Beardon~\cite{aB83} and Solynin~\cite{aS99,aS98}, 
 among
others.  One frequently cited open problem is to find 
for hyperbolically convex functions a sharp bound
on the Schwarz norm 
\[
\norm{S_f}_\D=\sup\{\abs{S_f(z)} \eta_\D^{-2}(z):\: z\in \D\},
\]
where $\eta_\D(z)=\frac{1}{1-\abs{z}^2}$
 is the hyperbolic density of $\D$ and $S_f$ is the Schwarzian derivative
\[
S_f = \left(\frac{f^{\prime\prime}}{f^\prime}\right)^\prime - 
\frac{1}{2}\left(\frac{f^{\prime\prime}}{f^\prime}\right)^2.
\]


The Schwarz norm
of an analytic function $f$ has long been a primary tool in understanding
its geometric behavior.  
For example, 
$\norm{S_f}_\D =0$ if and only if $f$ is a M\"obius transformation.
Thus $\norm{S_f}_\D$ is thought of as measuring how closely the geometric
behavior of $f$ resembles that of a M\"obius transformation.
Since the image of $\D$ under a M\"obius transformation must be a disc,
$\norm{S_f}_\D$ also measures the difference between the conformal geometry
of $f(\D)$ and that of a disc.  Lehto has used this idea to great effect,
producing a pseudo-metric on the set of all simply connected proper 
subdomains of $\C$.
See Lehto's book~\cite{oL87}, for example, for
a detailed discussion.  

 Nehari showed that if $f(\D)$ is convex (in the euclidean sense),
then $\norm{S_f}_\D\leq 2$, with equality if and only if $f(\D)$ is an 
``infinite strip'' bounded by two parallel lines~\cite{zN76}.
Similarly, Mej\'ia and Pommerenke showed that the extremal domain for spherically
convex functions is a ``spherical strip''~\cite{MP00b}.

The problem of finding a similar bound for the Schwarz norm of hyperbolically
convex functions has been intensely studied by a number of authors, including
Ma,  Minda,  Mej\'ia,  Pommerenke and  Vasilev~\cite
{MM94,MM99,MP00,MPV01,MP02,PM00}.  Mej\'ia and Pommerenke~\cite{PM00}
found partial results on the bound and conjectured that the
extremal value of $\norm{S_f}_\D$ is attained by a map of the form
\begin{equation}  \label{E:tanh}
f_\alpha (z)=\tan\left( \alpha\int_0^z \left(1-2\xi^2\cos 2\theta+\xi^4
\right)^{-1/2}\,d\xi \right),
\end{equation}
where $\alpha = \frac{\pi}{2K(\cos \theta)}$, 
and $K$ is the elliptic integral of the first kind.
The range of $f_\alpha$ is a ``hyperbolic strip''
bounded by two geodesics through $\pm i\tanh 
\left(\frac{\pi K(\sin\theta)}{4K(\cos\theta)}\right)$ and perpendicular
to the imaginary axis. 
See Figure~\ref{F:convex3}.


\begin{figure}
\begin{center}\scalebox{0.45}{
\epsfig{file=convex3} }
\caption{The extremal domains for maximizing the Schwarz norm in euclidean (left), spherical (center), and hyperbolic (right) geometry.}
\label{F:convex3}
\end{center}
\end{figure}


In this paper, we 
prove the following theorem, thus 
verifying Mej\'ia and Pommerenke's conjecture and  completing 
the classification of the extremal domains for the Schwarzian in all 
three of the classical geometries.


\begin{thm*}
The maximal value of the Schwarz norm for hyperbolically convex
functions is $S_{f_\alpha}(0)$, where
\[
f_\alpha(z) =\tan\left( \alpha\int_0^z \left(1-2\xi^2\cos 2\theta+\xi^4
\right)^{-1/2}\,d\xi \right),  \qquad \alpha  = \frac{\pi}{2K(\cos \theta)},
\]
$K$ is the elliptic integral of the first kind, and $\alpha$ is chosen
so that $\cos \theta$ is the unique critical point of the function
\[
g(s) = 4s^2-2 +\frac{\pi^2}{2K^2({s})}
\]
on $(0,1)$.
\end{thm*}

A computer calculation produces a  maximal value for the 
Schwarz norm for a hyperbolically convex
function of approximately $2.383635$. 
%\begin{thm}
%The maximal value of the Schwarz norm for a hyperbolically convex
%function is approximately $2.383635$ and is achieved by the ``hyperbolic strip
%map'' $f_\alpha$ with $\alpha\approx 0.5343598$.
%\end{thm}

\begin{remark}
Our variational techniques are used to argue that if an extremal
function exists, it must have certain properties.  
To guarantee the existence of an extremal function, we restrict our attention
to a dense collection of compact classes (hyperbolically convex
functions which map onto domains with at most a fixed number
of (proper) sides) and argue that an extremal function exists for these
classes.
Our methods show that the map $f_\alpha$ 
defined in~\eqref{E:tanh} 
is extremal in each such class for a common value of $\alpha$
that is approximately equal to $0.5343598$.
Moreover, since the classes are dense in the space of all hyperbolically 
convex functions, this $f_\alpha$ must be extremal for all hyperbolically
convex functions.  
Indeed, up to disc automorphisms, $f_\alpha$ can be the only extremal 
function which maps onto a finite-sided hyperbolically convex polygon.
However, this does not rule out the possibility, that some
function onto a non-polygonal region could have the same Schwarz
norm as $f_\alpha$. Thus our methods do not guarantee uniqueness,
but we conjecture that the extremal is in fact unique.
\end{remark}

%\begin{remark}

%Our arguments and the variations we construct can also be extended to  the 
%euclidean and spherical settings.  Thus we have a unified approach
%to all three classical geometries.
%\end{remark}

In Section~\ref{S:hypercon} we develop background material on
hyperbolic convexity and the Schwarzian derivative and discuss more of
the history of the problem.  In Section~\ref{S:prelim}
we first show by an approximation 
argument that the problem can be reduced
to studying the value of the Schwarzian at the origin for those
functions which map onto finite-sided polygons.
% to proving that 
%among all hyperbolically convex polygons with at most a fixed number
%of sides, the extremal domain has at most $2$ sides.  
%Thus  we need 
%only work with compact families where the existence of an extremal
%function is assured and our variational methods can be applied.  
%The $\sup$ can be replaced by $\max$ and


We then give a generalization of the Julia Variational Formula that
we will use in Sections~\ref{S:twovar} and~\ref{S:effect} to develop
two variations which preserve hyperbolic convexity. 
In Section~\ref{S:2sides} we  apply these variations to show the 
 extremal polygon has at most two proper sides.
Finally, in Section~\ref{S:compute} we employ inequalities on 
elliptic integrals to verify that the
two-sided domain conjectured by Mej\'ia and
Pommerenke~\cite{MP00,MP01,PM00} is indeed extremal.

We would like to thank the Phillip Brown, the referee and the editor for their 
very valuable suggestions.

%After developing some background material on hyperbolic convexity and a
%generalization of 
%the Julia variation in Sections~\ref{S:hypercon} and~\ref{S:julia}, 
%we develop two variations in Sections~\ref{S:twovar} and~\ref{S:effect}
%which preserve
%hyperbolic convexity. These variations 
%allow us to show in Section~\ref{S:2sides} that the
%extremal value occurs for a function whose range has at most two
%sides.  


\section{Hyperbolic Convexity and Schwarzians}\label{S:hypercon}

\subsection{Hyperbolic Geometry}

The unit disc $\D$ equipped with the metric
\[
d_h (z,w)=\inf
\left\{\int_\gamma \frac{1}{1-\abs{z}^2}\,\abs{dz} : \gamma 
\text{ is a rectifiable curve joining $z$ and $w$}\right\}
\]
forms a model for the hyperbolic plane~\cite{aB83}.  
Notice that the Poincar\'e
density 
\[
\eta_\D(z)=\frac{1}{1-\abs{z}^2}
\]
goes to infinity as $z$ moves toward the boundary of the disc.
Consequently, integrating $\eta_\D$ over curves near the boundary
produces large values of the integral.  If $z$ and $w$ do not lie on
a ray through the origin, then the euclidean line segment
joining them will produce a larger integral than a curve which bends 
away from the boundary.  In fact, the infimum  will be achieved by
an arc of a circle perpendicular to $\partial \D$.  Such curves are
hyperbolic geodesics.
Since disc automorphisms  
\[
M(z)=e^{i\theta}\frac{z-a}{1-\overline{a}z},
\]
where $\theta\in[0,2\pi)$ and $a\in\D$,
preserve circles 
orthogonal to $\partial \D$,
they are precisely the isometries of $\D$.

Any region $G$ conformally equivalent to $\D$ also carries a hyperbolic metric
defined in the same manner using  the density
\[
\eta_G(z)=\frac{\abs{f'(z)}}{1-\abs{f(z)}^2},
\]
where $f$ is a conformal map of $G$ onto $\D$.  Notice that it doesn't
matter which map $f$ is chosen as any two such maps must differ by a
disc automorphism.


\subsection{Convexity}

The euclidean notion of convexity generalizes to hyperbolic regions in an
obvious manner.

\begin{definition}A region $\Omega\subset\D$ is \textbf{hyperbolically convex}
if for any two points $z,w\in\Omega$, the hyperbolic  geodesic segment joining 
$z$ and $w$ lies completely in $\Omega$.
\end{definition}


Notice that since the disc automorphisms are the isometries
of the hyperbolic plane, the image $M(\Omega)$ of $\Omega$ under a disc
automorphism $M$ is hyperbolically convex if and only if $\Omega$ is.
The fundamental domains of discrete groups of disc automorphisms
provide a great many useful examples of hyperbolically convex domains.
See Beardon~\cite{aB83} for an extensive discussion of these regions.


We will call a
hyperbolically convex region $\Omega$ bounded by a finite number of
either geodesic arcs lying inside $\D$ or arcs of $\partial \D$ a
\textbf{hyperbolically convex polygon}.
We call the bounding geodesic arcs \textbf{proper sides} and the arcs
of $\partial \D$ \textbf{improper sides}.
For $n\geq 0$, we let 
\begin{align}\notag
K_n=\{&\text{hyperbolically convex polygons  containing $0$}\\\notag
& \text{and  having  
at most $n$ proper sides}\}
\cup\{0\}.
\end{align}
%Notice that $K_n$ contains 



\begin{definition}A conformal map $f:\D\to\Omega$ is called a
\textbf{hyperbolically convex function} if its range is hyperbolically convex.
We let $\mathbf {\mathit H}$ 
denote the class of all hyperbolically convex functions that fix the origin 
 and let $\bf{\mathit H_{\mathit n}}$ denote the subset of functions whose range
is in $K_n$.
\end{definition}

%A sequence of polygons in $K_n$ can converge (in the sense of Cartheodory) only to another polygon %in $K_n$. This convergence carries over to hyperbolically convex functions so that $H_n$ is
%compact.  On the other hand, every hyperbolically convex region is the limit of a sequence 
%of hyperbolically convex polygons; hence $\cup_n H_n$ is dense in $H$.  




\subsection{Schwarzians}

Much of the geometric behavior of an analytic function is described by
its Schwarzian derivative~\cite{GL2000,oL87}.

\begin{definition}
The \textbf{Schwarzian derivative} (or just ``Schwarzian'') of an
analytic function $f$ is
\[
S_f = \left( \frac{f''}{f'} \right)' - \frac12 \left( \frac{f''}{f'}
\right)^2 .
\].
\end{definition}


\begin{proposition} The Schwarzian of an analytic function is identically $0$
 if and only if it is a M\"obius 
transformation.  Moreover, the Schwarzian satisfies the chain rule
\[
S_{f\circ g}=\left(S_f \circ g\right)\left(g'\right)^2 + S_g.
\]
\end{proposition}

Thus, if $M$ is M\"obius, then 
\[
S_{M\circ g} = S_g.
\]
and
\[
S_{f\circ M}=\left(S_f \circ M\right)\left(M'\right)^2.
\]

Hence the Schwarzian is unchanged by post-composition with a M\"obius
transformation, but pre-composition produces an extra quadratic factor.

\begin{definition}
Let $f$ be defined on a simply connected region $G\subsetneq\C$.
The \textbf{Schwarz norm} of $f$ is given by
\[
\norm{S_f}_{G} = \sup_{z \in G} \eta_G^{-2}(z)\abs{S_f (z)}.
\]
\end{definition}

By taking into account the density of the hyperbolic metric, the
Schwarz norm is completely M\"obius invariant.  It is easy to show
for any M\"obius $M$ that 
\[
\eta_{M^{-1}(G)}(z)=\eta_{G} (M(z)) \abs{M'(z)},
\]
and thus
\[
\frac{\abs{S_{f\circ M}(z)}}{\eta^2_{M^{-1}(G)}(z)} =
\frac{\abs{S_f(M(z))}\,\abs{M'(z)}^2}{\eta^2_{G}(M(z))\,\abs{M'(z)}^2} =
\frac{\abs{S_f(w)}}{\eta^2_{G}(w)}.
\]
where $w=M(z)$.  Thus 
\[
\norm{S_f}_G=\norm{S_{f\circ M}}_{M^{-1}(G)}
\]
and
\[
\norm{S_f}_G=\norm{S_{M\circ f}}_G.
\]
In particular, notice that $\norm{S_f}_\D$ is unchanged by disc
automorphisms.



\subsection{Computing the Schwarz Norm of a Function}

%When computing the Schwarz norm of a function, 
%it can sometimes be very helpful to 
%change its domain.
Because of the M\"obius invariance of the Schwarz norm, we can 
change the domain of a function 
so that 
%take our function to be defined on a region where either the 
the hyperbolic density or the expression of the function itself is simpler.

\begin{example}\label{E:one-side}
Consider the function
\[
g(z)=\frac{\sqrt{z}-1}{\sqrt{z}+1}
\]
from the upper half plane $\HH$ onto the upper half-disc 
$\D^+=\{z\in\D:\:\im z>0\}$.  Notice that any hyperbolically convex 
domain bounded by exactly one proper side can be mapped onto $\D^+$ by
a disc automorphism.  Moreover, $g$ could have been defined on $\D$ by
precomposition with a M\"obius transformation; however, that would not
change the Schwarz norm.  Consequently, the Schwarz norm for any hyperbolically
convex function whose image is bounded by exactly one proper side must be equal to 
$\norm{S_g}_\HH$.

Next notice that $g$ itself is the composition of $h(z)=\sqrt{z}$ and another
M\"obius transformation.  Thus $\norm{S_g}_\HH=\norm{S_h}_\HH$.  It is 
now easy to compute
\[
S_h(z)=\frac{3}{8z^2}.
\]

The hyperbolic density on $\HH$ has an especially useful representation 
\[
\eta_\HH(z)=\frac{1}{2\im z}.
\]
Since both $\eta_\HH$ and $S_h$ are invariant under translations of the form $x + iy \rightarrow (x+a) + iy$, 
\[
\norm{S_h}_\HH=\sup_{z\in\HH} \eta^{-2}(z)\abs{S_h(z)} 
=\sup_{y>0} 4y^2\left|\frac{3}{8(iy)^2}\right| = \frac{3}{2}.
\]


\end{example}

%M\"obius invariance also allows another useful trick for computing the
%Schwarz norm. 



\subsection{Geometry of the Schwarzian}

Since $\Snf=0$ if and only if $f$ is M\"obius, we can view 
$\Snf$ as measuring how close $f$ is to being a M\"obius transformation.
Since any M\"obius transformation would send $\D$ to another disc or half plane, 
$\Snf$ also measures the amount of deformation between $f(\D)$
and a disc.  This notion
was formalized by Lehto~\cite{oL87} 
to produce a pseudometric between regions conformally equivalent to a disc.

There are a number of results that show that if $\norm{S_f}_\D$ is small,
then $f(\D)$ possesses disc-like properties.  The two most important
for our purposes are due to Nehari~\cite{zN49,zN76}.

\begin{theorem}  If $\Snf<2$, then $f$ is univalent and $f(\D)$ is a quasidisc.
Moreover, if $f$ is univalent, then $\Snf\leq 6$.
\end{theorem}

\begin{theorem} If $f(\D)$ is convex (in the euclidean sense),
 then $\Snf\leq 2$, with equality if and
only if $f(\D)$ is an infinite strip.
\end{theorem}

Mej\'ia and Pommerenke~\cite{MP00b} proved a similar result for spherically 
convex regions.

\begin{theorem} 
If $f(\D)$ is spherically convex, then 
\[
\norm{S_f}_\D \leq 2(1-\sigma(f)^2),
\]
where 
\[
\sigma(f)=\max_{z\in\D} (1-\abs{z}^2)\frac{\abs{f^\prime}}{1+\abs{f}^2}.
\]
For a fixed value of $\sigma(f)$, this
 maximum value of $\Snf$ is achieved by a map of the form
$f_\phi(z)=i\tanh\left(\frac{2\phi}{\pi}\Log
\left(\frac{1+z}{1-z}\right)\right)$ which takes $\D$ onto a 
 ``spherical strip,'' that is, a lune bounded by great circles through
$\pm i$ and making an angle $2\phi$ with the imaginary axis.
\end{theorem}

Thus, convex and spherically convex regions cannot be deformed
 too far from being a disc
in the sense of the Lehto pseudometric, and the regions 
with the greatest amount of deformation
are strips.  It has been conjectured by Mej\'ia and 
Pommerenke~\cite{MP00,PM00,MP01}
%and by Ma and Minda~\cite{MM99} 
that the same must hold for hyperbolically
convex regions.  


In particular, Mej\'ia and Pommerenke conjectured that the maximum value
of $\Snf$ over all hyperbolically convex maps $f$ is achieved by 
\begin{equation}\label{E:f_alpha}
f_\alpha(z) =\tan\left( \alpha\int_0^z \left(1-2\xi^2\cos 2\theta+\xi^4
\right)^{-1/2}\,d\xi \right),  \qquad \alpha  = \frac{\pi}{2K(\cos \theta)},
\end{equation}
where $K$ is the elliptic integral of the first kind.
The range of $f_\alpha$ is a ``hyperbolic strip''
bounded by two geodesics through $\pm i \tanh 
\left(\frac{\pi K(\sin\theta)}{4K(\cos\theta)}\right)$.
See Figure~\ref{F:strip}.

\begin{figure}
\begin{center}\scalebox{0.65}{
\epsfig{file=f_alpha-onto-2-sides} }
\caption{The hyperbolically convex extremal function for maximizing
  the Schwarz norm.  The extremal domain consists of an odd symmetric  polygon bounded by two proper geodesic sides - a ``hyperbolic strip.''}
\label{F:strip}
\end{center}
\end{figure}

%As our primary tool for proving their conjecture, we will
%modify  an extension of the  Julia variation.
%~\cite{BS84,BL75}.


\section{Preliminary Simplifications}\label{S:prelim}

\subsection{Reduction to Polygonal Domains}

As a first step toward verifying Mej\'ia and Pommerenke's conjecture,
we first reduce the problem to  hyperbolically convex polygons.
 
\begin{lemma} To determine the extremal value of $\norm{S_f}_\D$ over $H$,
it suffices to determine the extremal value over each $H_n$.
Moreover, each $H_n$ is compact.
\end{lemma}

\begin{proof}

From its definition as a supremum, it follows that
the Schwarz norm is a lower semi-continuous functional on $H$; 
that is, for any point $f \in H$, if $\{f_n\}$ is a sequence in
 $H$ which converges to $f$ (locally uniformly on compacta), 
then $\liminf \norm{S_{f_n}}_\D \ge \Snf$.  
To see this, let $\e > 0$.  
Then, for any compact subset $K$ of $\D$,
 there exists an index $N$ such that for $n \ge N$ we have
$\norm{S_{f_n\vert_{\interior K}}}_\D 
\ge \norm{S_{f\vert_{\interior K}}}_\D - \e$. 
 Hence, for $n \ge N$ we have 
$\norm{S_{f_n}}_\D\ge \norm{S_{f_n\vert_{\interior K}}}_\D 
\ge \norm{S_{f\vert_{\interior K}}}_\D - \e$.  
Hence, $\liminf \norm{S_{f_n}}_\D \ge \norm{S_{f\vert_{\interior K}}}_\D - \e$. 
 Since $K$ was abritrary we must have, $\liminf \norm{S_{f_n}}_\D \ge \Snf -
 \e$. 
 Since $\e$ was arbitrary, we have $\liminf \norm{S_{f_n}}_\D \ge \Snf$.  


For a fixed $n$, a sequence $D_k$ of hyperbolically convex 
polygonal domains
with at most $n$ (proper) sides can converge in the sense of 
Carath{\'e}odory~\cite{cP92}
only to another polygonal domain $D_0$, 
where $D_0$ has at most $n$ (proper) sides.  This 
follows because if $D_0$ had more than $n$ proper sides, 
then (from the boundary behavior) 
the Carath{\'e}odory  convergence~\cite{cP92}
would imply, for $k$ sufficiently large, 
that $D_k$ would have at least $n+1$ proper sides.  
From Carath{\'e}odory 's Convergence Theorem, 
this carries over to hyperbolically convex functions in $H_n$.  
Thus, $H_n$ is compact.

On the other hand, the boundary $\Gamma$ of a hyperbolically convex 
domain $D$ can be approximated by geodesics lying 
in $\overline{\D}$ by connecting points on $\Gamma$ 
sufficiently close together with geodesics lying 
in $\overline{\D}$.  It follows that we can find a sequence of hyperbolically
convex polygons that converge, in the sense of Carath\'eodory, to $D$.
Hence, % (by~\cite[cP92])
the mapping functions onto these polygons converge locally uniformly
to the map onto $D$. As a result, $\cup_n H_n$ is dense in $H$.  


Thus any extremal function in $H$ can be approximated by functions in $H_n$. 
As a result, it suffices to determine the supremum of 
$\norm{S_{f}}_\D$ in each $H_n$.

\end{proof}

%Consequently, we need only show that $f_\alpha$ is extremal in each $H_n$.  


\subsection{Reduction to $S_f(0)$}

Recall that the Schwarz norm was defined by taking the supremum 
of $\abs{S_f(z}\left(1-\abs{z}^2\right)^2$
over all points $z\in\D$.  Since 
$\abs{S_f(z}\left(1-\abs{z}^2\right)^2$ is invariant under disc automorphisms,
we need only consider $\abs{S_f(0)}$.

\begin{lemma}\label{L:sf0}
For each $n\geq 2$,
\[
\sup_{f\in H_n} \norm{S_f}_\D = \max_{f\in H_n} \abs{S_f(0)}.
\] 
\end{lemma}

\begin{proof}
Recall that every $f\in H_n$ is conjugate by a M\"obius transformation to
a map $g:\HH\to\HH$ which is hyperbolically convex in the upper half plane
model.  
Suppose  $g$ maps $\HH$ onto a curvilinear polygon with angles
$\alpha_1\pi,\alpha_2\pi,\dots, \alpha_n\pi$,  and the preimages of the
vertices are $a_1,a_2,\dots,a_n\in\R$. 
Then the Schwarzian of $g$ has the following very useful representation
(see Nehari~\cite{zN52})
\begin{equation}
S_g(z) = \frac{1}{2}\sum_{k=1}^n \frac{1-\alpha_k^2}{(z-a_k)^2} 
+\sum_{k=1}^n\frac{\beta_k}{z-a_k},
\end{equation}
where $\beta_1,\beta_2,\dots,\beta_n$ are real constants determined by
the angles.
Thus
\[
\norm{S_g}_\HH = \sup_{z\in\HH} 4(\im z)^2 
\left|\frac{1}{2}\sum_{k=1}^n \frac{1-\alpha_k^2}{(z-a_k)^2} 
+\sum_{k=1}^n\frac{\beta_k}{z-a_k}\right|.
\]

Notice however, that
\[
\lim_{\overset{z\to a\in \R}{a\neq a_k}} 4(\im z)^2 
\left|\frac{1}{2}\sum_{k=1}^n \frac{1-\alpha_k^2}{(z-a_k)^2} 
+\sum_{k=1}^n\frac{\beta_k}{z-a_k}\right| = 0
\]
and 
\[
\lim_{z\to a_k} 4(\im z)^2 
\left|\frac{1}{2}\sum_{k=1}^n \frac{1-\alpha_k^2}{(z-a_k)^2} 
+\sum_{k=1}^n\frac{\beta_k}{z-a_k}\right|
= 2(1-\alpha_k^2)<2.
\]
Applying the relationships between $\alpha_i$ and 
$\beta_i$ given by Nehari~\cite{zN52}, we can similarly
show 
\[
\lim_{z\to \infty} 4(\im z)^2 
\left|\frac{1}{2}\sum_{k=1}^n \frac{1-\alpha_k^2}{(z-a_k)^2} 
+\sum_{k=1}^n\frac{\beta_k}{z-a_k}\right|
\]
is also bounded by 2.

In view of the invariance of the Schwarz norm under
M\"obius transformations, the same bounds must be true for mappings of $\D$
onto hyperbolic polygons in $\D$. 
It is not difficult~\cite{PM00} to construct mappings of $\D$ onto 
hyperbolic polygons with $\norm{S_f}_\D \geq 2$.  Thus if 
$f$ maps $\D$ onto a hyperbolic polygon and $f$ is extremal for
maximizing the Schwarz norm, then the
supremum  in the definition of the Schwarz norm must in fact be attained at some
interior point.  
 
Suppose $f$ is extremal in $H_n$ for maximizing the Schwarz norm and
\[
\max_{z\in\D} (1-\abs{z}^2)^2\abs{S_f(z)}= (1-\abs{z_0}^2)^2\abs{S_f(z_0)}.
\]
Since $(1-\abs{z}^2)^2\abs{S_f(z)}$ is 
invariant under M\"obius transformations,
precomposing $f$ with a disc automorphism which sends zero to $z_0$
will produce a map $f_0$  which has the same Schwarz norm as $f$ but for
which $\norm{S_{f_0}}= \abs{S_{f_0}(0)}$.  

Notice by precomposing $f$ with a disc automorphism,
our new map $f_0$ need not fix $0$ and thus need not remain in $H_n$.
However, by postcomposing $f_0$ with a disc automorphism which 
sends $f_0(0)$ to $0$ and then rotates so that $S_{f_0}(0)$ is real,
  we can constuct an extremal function  $F\in H_n$ with
\[
\norm{S_f}_\D=\norm{S_F}_\D=S_{F}(0).
\]

Since the functional $S_f(0)$ is continuous on $H_n$ and $H_n$ is compact,
then
\[
\sup_{f\in H_n} \norm{S_f}_\D = \sup_{f\in H_n} S_f(0) = 
\max_{f\in H_n} S_f(0).
\]
\end{proof}







%Next we describe a variation we can apply to the class $H_n$.


\section{Variational Techniques}\label{S:twovar}

\subsection{The Julia Variation and Extensions}\label{S:julia}

Let $\Omega$ be a region bounded by a piecewise analytic curve $\Gamma$ and 
$\phi(w)$ be a positive piecewise $C^1$ function on $\Gamma$, vanishing
where $\Gamma$ is not analytic.  Denote the outward pointing 
unit normal vector at each point $w$ where $\Gamma$ is smooth by $n(w)$.
For $\epsilon$ near $0$, construct a new curve 
\[
\Gamma_\epsilon=\{w+\epsilon\phi(w)n(w):\: w\in\Gamma\}
\]
and let $\Omega_\epsilon$ be the new region bounded by $\Gamma_\e$.

For our purposes, we may assume that $\Omega$ is simply connected.  Thus
if $\epsilon$ is sufficiently small and
$\epsilon >0$, then $\Gamma$  is ``pushed outside'' the
domain, while if $\epsilon < 0$, then $\Gamma$  is ``pushed inside'' the
domain.
With the above notation, we state Julia's 
%modification of the
%Hadamard 
variational formula.
%; see~\cite{BL75,BS84} for similar variations of
%circular boundary arcs.

\begin{julia*}
Let  $f$ be a conformal map from $\D$ onto $\Omega$ with $f(0)=0$,
and suppose that $f$ has a continuous extension to $\partial\D$, which we also
denote by $f$. Then for $\e$ sufficiently small,
a similarly normalized conformal map $f_{\e}$ from $\D$ 
onto $\Omega_\epsilon$, 
with $f_{\e}(0) = 0$, is given by 

\begin{align}\label{E:julia}
f_\epsilon(z) = f(z)+\frac{\epsilon z f'(z)}{2\pi}\int_{\partial \D}
\frac{1+\xi z}{1-\xi z}\,d\Psi + E(\e,z),
\end{align}
where $d\Psi=\frac{\phi(f(\xi))}{\abs{f'(\xi)}}\,d\theta$ and $\xi=e^{i\theta}$
and $E(\e,z)$ is $o(\e)$ for $z$ on compact subsets of $\D$ and is continuously 
differentiable in $\epsilon$ for each fixed $z\in\D$.
\end{julia*}

Notice the restriction that $\phi$ vanish at the points of non-smoothness
of $\partial \Omega$ was a strong one.  It implies, for example, that while we
could vary the (proper) sides of a hyperbolic polygon, there is a difficulty
at the vertices.
However, it follows from the work of the first author
and J. Lewis~\cite{BL75} that 
such an extended version of the Julia 
 variation is possible (except at internal cusps, ie., 
where two proper sides meet at an angle 
of measure $2\pi$) which allows $f^\prime$ to vanish at the corners.
Moreover, they showed that the resulting function will agree
with the Julia variational formula on 
compact subsets up to $o(\epsilon)$ terms.
This will allow us to create variations which
 avoid any problems at the corners and preserve hyperbolic convexity.
 
%and apply the Julia formula to vary the (proper) 
%sides of hyperbolically
%convex polygons while preserving  the hyperbolic convexity.

\begin{figure}
\begin{center}\scalebox{0.7}{
\epsfig{file=julia-out} }
\caption{The variation produced by ``pushing out'' a side.}
\label{F:julia-out}
\end{center}
\end{figure}

\subsection{Two Variations}\label{SS:twovar}
We next describe in detail two variations for functions 
in $H_n$ which preserve hyperbolic convexity.
First, if $f$ maps onto a hyperbolically
 convex polygon $\Omega$ with a proper side $\Gamma$,
we can ``push'' $\Gamma$ to a nearby geodesic  $\Gamma_\epsilon$ 
in such a way that the varied
function $f_\epsilon$ will still be hyperbolically convex.  See
Figure~\ref{F:julia-out}.



\begin{lemma} \label{L:julia-out}
Suppose $f\in H_n$ and $f$ is not constant nor the identity map.
If $\Gamma=f(\gamma)$ is a proper side of $\Omega=f(\D)$, 
then for $\epsilon$ sufficiently small
there exists a variation $f_\epsilon\in H_n$ which ``pushes'' 
$\Gamma$ to a nearby geodesic $\Gamma_\epsilon$,  
% That is, $f_\epsilon(\D)$ 
%s a hyperbolic polygon whose sides are the same as those of $f$, except that
%$\Gamma$ has been replaced by $\Gamma_\epsilon$, 
where $\Gamma_\epsilon\to\Gamma$
as $\epsilon\to 0$. 
% is a geodesic through $m+\epsilon n(m)$ and $m$ is the euclidean midpoint
% of $\Gamma$.
This variation, $f_\epsilon$, agrees with the Julia Variational Formula 
on compact subsets up to
$o(\epsilon)$ terms.
\end{lemma}

\begin{proof}

Let $\Gamma$ be a (proper) 
side of the boundary of $\Omega$ which lies on a circle $\Lambda$.  
Without loss of generality we can assume that $\Omega$ 
has been rotated so that the center
of $\Lambda$ lies on the positive real-axis. To define our variation first map
$\Omega$ into the right half-plane under the map $h(z) = \frac{1+z}{1-z}$ and 
let $\Upsilon = h(\Gamma)$ and $\Delta  = h(\Omega)$. Let $g = h \circ
 f$.  We note that
$\Delta = g(\D)$.

Under $h$, the geodesics in 
$\D$ map to arcs in the right half-plane which lie on circles whose centers are 
on the imaginary axis. We note that $h(\Lambda)$ is a circle centered at the origin. 
Let the arc $\widehat{AB}$ denote $\Upsilon $ on the circle $h(\Lambda)$ and let 
$r_{\Lambda}$ be the radius of $h(\Lambda)$.  

Now, considering the case where $\e > 0$, i.e., the
case where we ``push out'' a (proper) side of the domain by a euclidean
distance $\e$ and thus  enlarge the domain.  Let the 
circle $\chi$ of radius $r = r_{\Lambda} - \e$ be obtained by
multiplying
 $h(\Lambda)$ by 
$1- \e/r_\Lambda$
 and let $\widehat{CD}$ be the arc of $\chi$ which connects the extensions 
of the circles defined by the images of the $\Gamma$'s adjacent geodesics under $h$, 
when $A$ and $B$ are in the right-half plane. If $A$ or $B$ is on the imaginary axis, 
then $C$ or $D$ will be the point on the imaginary axis where $\chi$ intersects the 
imaginary axis, resp.  In this case, the interval $[A,C]$ or $[D,B]$ will be added to 
the boundary of our varied domain, resp.

Now recall in the original Julia variation formula, the function
 $\phi$
 must be $0$ at 
the end points of $\widehat{AB}$.  Thus at each end of $\widehat{AB}$,
we construct approximating
smooth curves $\sigma_A$ and $\sigma_B$ which connect $\widehat{AB}$ to 
$\widehat{CD}$ such that $\sigma_A$ and $\sigma_B$ are orthogonal to $\widehat{AB}$ 
at the endpoints $A$ and $B$ and $\sigma_A$ and $\sigma_B$ intersect 
$\widehat{CD}$ tangentially at points $C'$ and $D'$, resp.  They can be chosen to be 
within $\e^2$ of the arc they are approximating.  See Figure~\ref{F:julia-half-plane}.


\begin{figure}
\begin{center}\scalebox{0.7}{
\epsfig{file=julia-half-plane} }
\caption{The geodesics in the right half plane produced by our variation.}
\label{F:julia-half-plane}
\end{center}
\end{figure}


Now for $w \in \widehat{AB}$ let $n(w)$ be the outward 
unit normal to $g(\Lambda)$ and define 
$\phi(w) = |w^*-w|$ where $w^*$ is the nearest point on 
$\sigma_A \cup \widehat{C'D'} \cup \sigma_B$
along the normal $n(w)$.  Let the domain $\Delta^*_{\e}$ 
be the variation of $\Delta$ where the
arc $\widehat{AB}$ on the boundary of $\Delta$ has been 
replaced by 
$\sigma_A \cup \widehat{C'D'} \cup \sigma_B$. To obtain a variation 
of $\Delta$ which preserves 
hyperbolic convexity, 
let the domain $\Delta_{\e}$ be the
variation of $\Delta$ where the arc $\widehat{AB}$ has been replaced by 
$\widehat{CD}$ and 
connecting intervals, if necessary.  Let $g^*_{\e}$ and $g_{\e}$
be the corresponding 
mapping functions
onto $\Delta_\e^*$ and $\Delta_\e$, resp.  Let $\phi(w) = 0$ for 
$w \in \partial \Delta \setminus \widehat{AB}$.
Then, by the Julia Variational Formula, $g^*_{\e}$ satisfies

\begin{equation*}
g^*_{\e}(z) = g(z)+\frac{\e z g'(z)}{2\pi}\int_{\partial \D}
\frac{1+\xi z}{1-\xi z}\,d\Psi + o(\e),
\end{equation*}

\noindent on compact subsets, 
where $d\Psi=\frac{\phi(g(\xi))}{\abs{g'(\xi)}}\,d\theta$ 
and $\xi=e^{i\theta}$. Rewriting $g$ as
$h \circ f$, we have

\begin{equation*}
g^*_{\e}(z) = h(f(z))+\frac{\e z h'(f(z))f'(z)}{2\pi}\int_{\partial \D}
\frac{1+\xi z}{1-\xi z}\,d\Psi + o(\e),
\end{equation*}

\noindent on compact subsets, 
where $d\Psi=\frac{\phi(h(f(\xi)))}
{\abs{h'f((\xi))f'(\xi)}}\,d\theta$ and $\xi=e^{i\theta}$. 

The function $g^*_{\e}$ was constructed so that it satisfies the 
requirements for representation by 
the Julia Variational Formula.  The function $g_{\e}$ was constructed 
so that image, $g_{\e}(\D)$, was 
hyperbolically convex.  The arguments used in \cite{BL75}, pp 348-356, 
show that $g^*_{\e} - g_{\e} = o(\e)$ on compact subsets.  Then, defining 
$f_{\e} = h^{-1} \circ g_{\e} \equiv \frac{g_{\e}-1}{g_{\e}+1}$, 
a straightforward 
argument shows that 

\begin{equation*}
f_\epsilon(z) = f(z)+\frac{\epsilon z f'(z)}{2\pi}\int_{\partial \D}
\frac{1+\xi z}{1-\xi z}\,d\Psi + o(\e)
\end{equation*}
on compact subsets, 
where $\Psi$ is a real measure on $\partial \D$ and $\xi=e^{i\theta}$.  
Thus, $f_{\e}$ maps $\D$
onto a hyperbolically convex domain, i.e., $f_{\e} \in H_n$ for 
$\e$ sufficiently small.

A similar variation can be defined for $\e < 0$ that ``pushes in'' the 
(proper) side $\Gamma$ to a nearby geodesic.

\end{proof}



%Applying this technique on a subarc of $\Gamma$ ending on an endpoint of
%$\Gamma$,
If $\Gamma$ intersects another side $\Gamma^*$ at $z^*$, 
 and $z_0\in\Gamma$,  
then we can vary $f$ so as to replace
the portion of $\Gamma$ between $z_0$ and $z^*$ with a new geodesic between
$z_0$ and some $z^*_\epsilon\in\Gamma^*$.  See Figure~\ref{F:kick-in}.
A similar argument to the above then gives: 


\begin{lemma} \label{L:kick-in}
Suppose $f\in H_n$ and $f$ is not constant nor the identity map,
and $\Gamma=f(\gamma)$ is a proper side of $f(\D)$ meeting a side $\Gamma^*$.
Then, there exists a variation $f_\epsilon\in H_{n+1}$ which 
adds a proper side to $f(\D)$ by pushing one end of $\Gamma$ to
 a nearby side $\Gamma_\epsilon$.  That is, $f_\epsilon(\D)$ 
is a hyperbolic polygon whose sides are the same as those of $f$, except that
one end of $\Gamma$ has been replaced by $\Gamma_\epsilon$ and $\Gamma^*$
has been shortened.
Moreover, $f_\epsilon$ agrees with the Julia Variational Formula on compact
subsets up to
$o(\epsilon)$ terms.
\end{lemma}


\begin{figure}
\begin{center}\scalebox{0.65}{
\epsfig{file=julia-kick-in} }
\caption{The variation produced by pushing in one end of a side.}
\label{F:kick-in}
\end{center}
\end{figure}

\begin{remark}\label{R:end}
Notice that in order to maintain hyperbolic convexity, we can only push in 
the ``end'' of a side, that is, a subarc that ends at a vertex of the
polygon.  However, we can choose this  subarc to be as long, or
more importantly, as short, as we wish.  
\end{remark}


\section{The Effect of the Variations on the Schwarzian}\label{S:effect}

%If the supremum in the definition of the Schwarz norm occurs at some
%point $z_0$ in $\D$, we could pre-compose $f$ with a disc automorphism
%sending $0$ to $z_0$.  This will leave both the Schwarz norm and the
%range of $f$ unchanged.
%for many purposes we can assume $f$ has been normalized so that
%$\Snf=S_f(0)$.  

Recall from Lemma~\ref{L:sf0}, we can assume our extremal function
satisfies $\norm{S_f}_\D = S_f(0)$.
In this case, the Schwarz norm has a simple representation 
in terms of power series coefficients~\cite{oL87,PM00}.

\begin{proposition}\label{P:Sf0}
If $f=\alpha(z + a_2 z^2 + a_3 z^3 +\dots)$, 
% is extremal over $H_n$ for $\norm{S_f}_\D$, 
then
% has been normalized so that the sup in the definition of the
%Schwarz norm occurs at $z=0$, then
\[
S_f(0)=6(a_3 - a_2^2).
\]
\end{proposition}

\begin{proof}

Since $f'(0)=\alpha$, $f''(0)={2}\alpha a_2$, and 
$f'''(0)={6}\alpha a_3$,
\begin{align}\notag
\frac{d}{dz} \frac{f''(z)}{f'(z)}\bigg\vert_{z=0}
&=\frac{f'(0)f'''(0)- (f''(0))^2}{(f'(0))^2}\\\notag
&=\frac{6\alpha^2 a_3 - {4}\alpha^2 a_2^2}{\alpha^2}\\\notag
&=6a_3 -4 a_2^2.
\end{align}

Also,
\[
\frac{1}{2}\left(\frac{f''(0)}{f'(0)}\right)^2 
=\frac{1}{2}\frac{4 \alpha^2 a_2^2}{\alpha^2}
= 2 a_2^2.
\]

Consequently,
\[
S_f(0)= %\eta_\D(0) 
\left(
\left(\frac{f''}{f'} \right)' - \frac12 \left( \frac{f''}{f'}
\right)^2\right) = 6(a_3-a_2^2).
\]

\end{proof}

Now applying the representation for the variations provided by the
Julia Formula, we can calculate to first order in $\epsilon$ the
Schwarzian of our varied functions.

\begin{lemma}\label{L:Sg}
If $f(z)=\alpha(z +a_2 z^2 +a_3 z^3 + \dots)\in H_n$ 
%has been normalized so that $\norm{S_f}_\D=S_f(0)$ 
and either of the variations described in Lemma \ref{L:julia-out} and Lemma \ref{L:kick-in} 
are applied to $f$, then
the new function $f_\epsilon(z)$ has $S_{f_\epsilon}(0)$ given by
\[
6\left( \left( a_3 +\frac{\epsilon}{2\pi} \int_\Gamma
(3a_3 + 4a_2\xi + 2\xi^2) \, d\Psi \right) -
\left( a_2 +\frac{\e}{2\pi} \int_\Gamma (2a_2 + 2\xi) \, d\Psi \right)^2\right)
+o(\epsilon).
\]
\end{lemma}

\begin{proof}

Let $\gamma$ be the preimage in $\partial \D$ of the portion of the 
boundary we are varying.  Expanding $\frac{1+\xi z}{1-\xi z}$ as a 
series, we obtain on compact subsets

\begin{align}\notag
f_\epsilon(z)&=f(z) + \frac{\epsilon z f^\prime(z)}{2\pi} 
\int_\gamma \frac{1 +\xi z}{1- \xi z} \,d\Psi + o(\epsilon) \\\notag
%&= f(z) + \frac{\epsilon z f^\prime(z)}{2\pi} 
%\int_\gamma (1+\xi z)(1+ \xi z + \xi^2 z^2 + \xi^3 z^3 +\dots) \, d\Psi 
%+ o(\epsilon)\\\notag
&= f(z) + \frac{\epsilon z f^\prime(z)}{2\pi} 
\int_\gamma (1+2 \xi z+2 \xi^2 z^2 +2 \xi^3 z^2 +\dots)\, d\Psi+o(\epsilon).
\end{align}

Now since $zf'(z)=\alpha(z+2a_2 z^2 + 3 a_3 z^3 +\dots)$, we have
\begin{align}\notag
f_\epsilon(z) = 
% &f(z) + \frac{\epsilon}{2\pi} 
% \int_\gamma \alpha(z+2a_2 z^2 + 3 a_3 z^3 +\dots) \\\notag
% &(1+2 \xi z+2 \xi^2 z^2 +2 \xi^3 z^2 +\dots)\, d\Psi+o(\epsilon)\\\notag = 
& f(z)+\frac{\epsilon\alpha}{2\pi}\int_\gamma 
[z + 2(a_2 + \xi) z^2 + (3a_3 + 4a_2 \xi +2\xi^2)z^3 +\dots] \, d\Psi +
o(\epsilon)
\end{align}
on compact subsets.

Finally, gathering the powers of $z$, we arrive at
\begin{align}\label{E:fepsilon}
f_\epsilon(z)=\alpha\Biggl(& 
\biggl(1+\frac{\epsilon}{2\pi} \int_\gamma \, d\Psi\biggr) z+
\biggl(a_2+\frac{\epsilon}{2\pi} \int_\gamma 2(a_2+\xi)\, d\Psi\biggr)z^2\\\notag
& +
\left(a_3+\frac{\epsilon}{2\pi} \int_\gamma (3a_3+4a_2\xi+2\xi^2) \, d\Psi\right)z^3
+ \dots \Biggr)+o(\epsilon)
\end{align}
on compact subsets.

Hence by Proposition~\ref{P:Sf0},
\begin{align}\notag
\frac{1}{6} S_{f_\epsilon}(0) = 
&\left(a_3 +\frac{\epsilon}{2\pi}\int_\gamma (3a_3+4a_2\xi+2\xi^2) 
\, d\Psi\right)-
\\\notag
&\left(a_2+ \frac{\epsilon}{2\pi}\int_\gamma 2(a_2+\xi)\,d\Psi\right)^2
+o(\epsilon).
\end{align}


\end{proof}

Taking the derivative of the new Schwarzian  with respect to $\epsilon$, we have
\begin{align}\notag
\frac{1}{6}\frac{\partial}{\partial \e} S_{f_\epsilon}(0) = 
&\frac 1{2\pi}\iG (3a_3 + 4a_2\xi + 2\xi^2) \, d\Psi  \\ \notag
 &- 2 a_2\frac{1}{2\pi}
\iG 2(a_2+\xi) \, d\Psi - \frac{2\e}{2\pi}\left(\iG 2(a_2 + z) \, 
d\Psi\right)^2  +\dots \\\notag
& +o(1).
\end{align}

When $\e=0$, the derivative becomes
\begin{align}\label{E:Sg}
\frac{\partial}{\partial \e} S_{f_\epsilon}(0) \bigg\vert_{\e=0} 
%= &6\left(
%\frac{1}{2\pi}\iG 3a_3 +4a_2\xi +\xi^2 -4a_2^2 - 4a_2\xi \, d\Psi\right) \\\notag
&=\frac{6}{2\pi}\iG (3a_3-4a_2^2 + 2\xi^2) \, d\Psi.
\end{align}




\section{Reduction to at Most Two Proper Sides}\label{S:2sides}
With these computations, we can now show that any function mapping onto 
a region with
more than two disjoint proper sides cannot be extremal.  We do this in two 
steps, first reducing the extremal domain to one having at
 most four proper sides, then employing our
Step Down Lemma to further reduce the possibilities to 
at most two proper sides.


\begin{lemma}\label{L:4sides}
Suppose $f \in H_n$ is a hyperbolic convex function onto a region with
more than four proper sides, then $f$ cannot be extremal for maximizing
the Schwarz norm.
\end{lemma}


\begin{proof}
Suppose $f\in H_n$ is extremal, that $\norm{S_f}_\D = S_f(0)$
 and that $f$ maps onto a region with more than
four proper sides.  
%In Lemma~\ref{L:sf0}  we showed
%that we may assume the extremal function $f$ satisfies  
%$\norm{S_f}_\D = S_f(0)$.
We will apply our variations to show that $f$
cannot be extremal.  

%Recall, however, that our variations were
%designed to control $\abs{S_f(0)}$, not $\norm{S_f}_\D$ itself.



%the supremum in the definition of the Schwarz norm for 
%the extremal $f$ is acutally a
%maximum, and the maximum is atained at some interior point $z_0\in\D$.
%Thus we can pre-compose $f$ with a disc automorphism
%which sends $0$ to $z_0$ and produce a function with the same range,
%but normalized so that
%\begin{equation}\label{E:normalized}
%\norm{S_f}_\D=S_f(0). 
%\end{equation}

%On the other hand, given a function onto a hyperbolically
%convex $n$-gon which has been normalized
%so that~\eqref{E:normalized} holds, pre- and post-composition
%with disc automormhisms can be used produce a map which fixes the
%origin and thus lies in $H_n$.  Thus we may work with either
%normalization
%and in light of our variational lemmas, 
%  we assume that

%If $f$ is extremal, then 
Recall from equation~\eqref{E:Sg} that for either of our variations along a side
$\Gamma=f(\gamma)$
\[
\frac{\partial}{\partial \epsilon} S_{f_\epsilon} (0) \bigg|_{\epsilon=0}
  = \frac6{2\pi} \iG \left(3a_3 - 4a_2^2 + 2\xi^2\right) d\Psi.
\]


By composing $f_\epsilon$ with a rotation (which will not change $S_{f_\epsilon}$),
 we may assume
\[
\frac{\partial}{\partial \epsilon} S_{f_\epsilon} (0) \bigg|_{\epsilon=0}
=\frac{6}{2\pi} \iG \re \{3a_3 - 4a_2^2 + 2\xi^2\} \,d\Psi.
\]


But now notice that $\K (\xi) = 3a_3 - 4a_2^2 + 2\xi^2$ wraps
$\{\abs{\xi}=1\}$
 twice around the image circle $C$.  If $\K(\gamma)$ lies
completely in the right half plane, then 
we push the (proper)
side $\Gamma$ out, with $\e > 0$ as in Lemma \ref{L:julia-out}.  Then,
since $\re \K(\gamma) > 0$ for $\gamma = f^{-1}(\Gamma)$, we will have
that $\iG \re\{3a_3 - 4a_2^2 +
2\xi^2\} \, d\Psi > 0 $.  Similarly, if $\K (\gamma)$ lies completely in the
left half plane, $\Gamma$ can be pushed in so that
$\iG \re\{3a_3 - 4a_2^2 + 2\xi^2 \} \, d\Psi < 0$.  

On the other hand, 
if $f$ is extremal, then $\frac{\partial }{\partial \epsilon} S_{f_\epsilon}(0)
|_{\epsilon=0}$ must equal $0$. Otherwise, 
\[
\norm{S_f}_\D=S_f(0)<S_{f_\epsilon}(0)\leq\norm{S_{f_\epsilon}}_\D
\]
for some $\epsilon$ sufficiently near $0$.
Thus
$\K(\xi)$ must cross the imaginary axis at least once for every proper side.

But as $\xi$ travels along $\partial \D$, $\K(\xi)$ crosses the imaginary
axis at most 4 times.  Consequently, the range of our extremal function
$f$ can have at most $4$ proper sides.
\end{proof}

\begin{remark}  We note that the technique employed thus far is fairly 
general and can be used to show that an extremal domain for
general extremal problems has at most a small number of proper sides. 
But usually the number of parameters is still too large to
completely determine the extremal domain.  The following critical
Step Down Lemma enables us to further reduce the number of proper sides 
to the point that the actual maps are computable.
\end{remark}

%\begin{step-down}
%Consider the problem of finding the extremal value of some
%functional $\Phi$ over a class of functions $C$ mapping onto
%finite sided polygons.
%Suppose that a polygon with more than $N$ sides cannot be 
%the range of an extremal function.  
%Suppose there exists a
% generalized Julia variation $f_\epsilon$ of $f$ varying any side
%(or portion of a side) $\Gamma=f(\gamma)$ 
%of the range of $f$ so that 
%\[
%\frac{\partial}{\partial\epsilon} \Phi(f_\epsilon) \bigg\vert_{\epsilon=0}
%= \int_\gamma \re \K(\xi)\,d\Psi
%\]
%where the curve $\K(\xi)$ crosses the imaginary axis $2N$ times.
%Then the extremal function can have at most $N$ sides.
%\end{step-down}



\begin{step-down*}~\label{L:2sides}
%Applying a  variation which adds a side to the range, 
A hyperbolically convex function $f\in H_n$  onto a region with
more than two proper sides cannot be extremal for maximizing 
the Schwarz norm.
\end{step-down*}


\begin{proof}

If $f\in H_n$ is extremal, then by Lemma~\ref{L:4sides}, $f(\D)$ can 
have at most four proper sides. Thus $f\in H_4\subsetneq H_5$.

 If $f(\D)$ has exactly four proper sides,
then an application of the pigeonhole principle implies that
one end point $\K(z^*)$
of one of the arcs $\K (\gamma)$ must lie in the left half plane.
Thus
for some $z_0$ sufficiently close to $z^*$, the image under $\K$ of the 
subarc $\gamma_0$ of $\gamma$ joining $z_0$ and $z^*$ lies completely in the 
left half plane.  As
a result,
\[
\int_{\gamma_0} \re \{3a_3 - 4a_2^2 + 2\xi^2 \} \,d\Psi < 0.
\]

Applying Lemma~\ref{L:kick-in} to push
 the corresponding arc $\Gamma_0=f(\gamma_0)$ ``in'' (varying by
 -$\epsilon$),
 will then produce both an additional proper side and 
an increase in the Schwarzian since $-\int_\gamma \K(\xi)\, d\Psi >0$.
  This varied function $f_\epsilon$ will
thus lie in $H_5\setminus H_4$ and have a larger Schwarz norm than any
function in $H_4$.  But this contradicts Lemma~\ref{L:4sides},
since the extremal function in $H_5$ can have at most
$4$ proper sides.
As a result, the extremal domain for a map in $H_5$ can have at most $3$ 
proper sides.

%Since $H_4\subset H_5$, 
%But 
%since the extremal value is monotonic as a function of $n$, this 
%would violate
%the fact that the extremal domain for $H_5$ has at most 4 sides.  
%the extremal domain for $H_4$ cannot have 4 sides.


However, if the extremal domain has exactly 3 proper
sides, some endpoint must again lie
in the left half-plane.  As before, we may add a proper side and
increase the Schwarz norm.  Thus the extremal domain can have at most
two proper sides.


\end{proof}





\section{Computing the Schwarz Norms for One- 
and Two-Sided Functions}\label{S:compute}

Now we need only compute the Schwarz norm for the remaining possible 
extremal functions to complete the proof of our theorem. First, however,
 we need a couple of facts about the behavior of 
elliptic integrals~[1,~pp.~53-54].

\begin{lemma}~\label{L:matti}
If $K$ and $E$ are the complete elliptic integrals of the first and 
second kind, respectively, then
\begin{enumerate}
\item For each $c\in[1/2,2]$, the function $g(x)=(\sqrt{1-x^2})^c K(x)$
is decreasing from $[0,1)$ onto $(0,\pi/2]$.
\item The function 
\[
h(t)=\frac{E(t)-(1-t^2)K(t)}{t^2}
\]
is increasing from $(0,1)$ onto $(\pi/4,1)$.
\end{enumerate}
\end{lemma}

With this lemma, we can prove our main result:

\begin{thm*}
The maximal value of the Schwarz norm for hyperbolically convex
functions is $S_{f_\alpha}(0)$, where
\[
f_\alpha(z) =\tan\left( \alpha\int_0^z \left(1-2\xi^2\cos 2\theta+\xi^4
\right)^{-1/2}\,d\xi \right),  \qquad \alpha  = \frac{\pi}{2K(\cos \theta)},
\]
$K$ is the elliptic integral of the first kind, and $\alpha$ is chosen
so that $\cos \theta$ is the unique critical point of the function
\[
g(s) = 4s^2-2 +\frac{\pi^2}{2K^2(s)}
\]
on $(0,1)$.
\end{thm*}


Numerical calculation indicates the extremal value of the 
Schwarz norm is approximately $2.383635$ and is achieved by
$S_{f_\alpha}(0)$ with $\alpha\approx 0.5343598$.


\begin{proof}

By Lemma~\ref{L:4sides} 
and the Step Down Lemma, the extremal function in $H_n$, $n>2$, cannot
map onto a region with more than $2$ proper sides.  Thus we need only
determine whether the extremal function has $0$, $1$, or $2$ proper sides.

The only hyperbolically convex functions onto a region with no proper sides
are disc automorphisms which have
Schwarz norm equal to $0$.

Any hyperbolically convex function onto a region with exactly one proper side
differs from the function $g$ of Example~\ref{E:one-side} by 
composition with M\"obius transformations,
and hence has the same Schwarz norm, namely $3/2$.  Note that 
Mej\'ia and Pommerenke~\cite{PM00} also explicitly computed these maps
and their Schwarz norm.


Next notice that for two-sided domains whose proper sides 
intersect inside the disc, the intersection point
can be moved to the origin by a disc automorphism.
Moreover, domains bounded by two proper sides which
do not intersect 
can be mapped via a disc automorphism onto an odd symmetric domain,
as in figure~\ref{F:strip},
and domains whose two proper sides intersect on the boundary of the disc 
can be approximated by domains whose (proper) sides do not  intersect.  
Thus there are two types of regions bounded by exactly two proper sides that we
must consider: 
\begin{enumerate}
\item \label{I:origin}  Domains 
in which the two proper sides intersect at the origin
\item Odd symmetric domains in which the proper sides do not intersect
\end{enumerate}

For case~\eqref{I:origin}, consider the family of maps
\[
g_\alpha(z)= \left(\frac{\sqrt{z}-1}{\sqrt{z}+1}\right)^\alpha,\quad
0<\alpha<1
\]
from the upper half plane onto a sector of $\D$ with an opening of 
angle $\alpha\pi$ at the origin.
Functions which map onto domains of the first type can be obtained from
one of the maps $g_\alpha$ by composition with M\"obius transformations, thus
we may limit our attention to this family.

Computation and simplification reveals
\[
S_{g_\alpha}(z)=\frac{3z^2-(2+4\alpha^2)z +3}{8z^2(z-1)^2}.
\]
Because both the Schwarzian and the hyperbolic density $\eta_\HH(z)=1/(2\im{z})$
are invariant under horizontal translation, we need only consider
pure imaginary numbers when computing the Schwarz norm.
Thus
\begin{align}\notag
\norm{S_{g_\alpha}}_\HH=&\sup_{z\in\HH} 4\left(\im z\right)^2
\left|\frac{3z^2-(2+4\alpha^2)z +3}{8z^2(z-1)^2}\right| \\\notag
&=\sup_{y>0} \frac{1}{2}\frac{\abs{3y^2+(2+4\alpha^2)iy -3}}{y^2+1}.
\end{align}

The derivative of the right hand side with respect to $y$ is
\[
\frac{-8(\alpha^2-1)(\alpha^2+2)(y^3-y)}
{\sqrt{9y^4+(16\alpha^4+16\alpha^2-14)y^2 +9}(y^2+1)^2}.
\]
Clearly the denominator is positive for all $y>0$, $0<\alpha<1$.
Moreover, for all $0<\alpha<1$ the numerator is negative for 
$0<y<1$ and positive for $1<y<\infty$.  
Consequently, $\eta_\HH(iy)\abs{S_{g_\alpha}(iy)}$ is decreasing 
on $0<y<1$ and increasing on $1<y<\infty$.
Thus the supremum in
the definition of $\norm{S_{g_\alpha}}_\HH$ must occur as $y$ approaches $0$
or as $y$ approaches $\infty$.  See Figure~\ref{F:g-alpha}.


\begin{figure}
\begin{center}\scalebox{0.45}{
%\begin{rotate}{-90}
%\begin{minipage}{\linewidth}
\epsfig{file=g-alpha,angle=-90} }
%\end{minipage}
%\end{rotate} }
\caption{The graph of $4y^2\abs{S_{g_\alpha}(iy)}$ for $\alpha=0.75$.}
\label{F:g-alpha}
\end{center}
\end{figure}

In both directions, however,
\begin{align}
\lim_{y\to 0^+}\eta_\HH(iy)\abs{S_{g_\alpha}(iy)} & = 
\lim{y\to 0^+} \frac{1}{2}\frac{\abs{3y^2+(2+4\alpha^2)iy -3}}{y^2+1} 
=\frac{3}{2}\\\notag
\lim_{y\to \infty}\eta_\HH(iy)\abs{S_{g_\alpha}}(iy) & = 
\lim{y\to \infty} \frac{1}{2}\frac{\abs{3y^2+(2+4\alpha^2)iy -3}}{y^2+1} 
=\frac{3}{2}.
\end{align}
Notice that both limits are independent of $\alpha$.  Consequently, 
the Schwarz norm of any hyperbolically convex function onto a domain
bounded by exactly two proper sides intersecting inside $\D$ must be $3/2$.

The second case consisting of maps onto odd symmetric two-sided domains
can be handled using the family of maps 
$f_\alpha$ constructed by Mej\'ia and Pommerenke~\cite{PM00}.
Each hyperbolically convex function
\[
f_\alpha(z)=\tan\left( \alpha\int_0^z \left(1-2\xi^2\cos 2\theta+\xi^4
\right)^{-1/2}\,d\xi \right),
\]
maps onto a ``hyperbolic strip''
bounded by two geodesics through 

$$
\pm i \tanh \left(\frac{\pi K(\sin\theta)}{4K(\cos\theta)}\right),
$$

\noindent where  $\alpha  = \frac{\pi}{2K(\cos \theta)}$, $0<\theta<\pi/2$,
and $K$ is the elliptic integral of the first kind.
See Figure~\ref{F:strip}.


For these functions $f_\alpha$, our problem is to compute
\[
\norm{S_{f_\alpha}}_\D = \sup_{z\in\D} \left(1 - \abs{z}^2\right)^2 
\abs{S_{f_\alpha}(z)}.
\]
Using an extensive computational argument which considers several cases
(various interval ranges for $\abs{z}$, $\arg z$, and $\alpha$) and
uses properties of polynomials and $K$, one can show that this 
problem can be reduced to computing
\[
\sup_{0\leq x<1}\left(1-x^2\right)^2 \abs{S_{f_\alpha}(x)}.
\]



Mej\'ia and Pommerenke~\cite{PM00} computed $S_{f_\alpha}(x)$ to be
\[
S_{f_\alpha}(x)=\frac{2( (\cos(2\theta)+\alpha^2) 
+ (\cos^2(2\theta)-2\alpha^2\cos(2 \theta) -3)x^2 
+ (\cos(2\theta) +\alpha^2)x^4)}{(1-2\cos(2\theta)x^2 +x^4)^2},
\]
where $\alpha=\frac{\pi}{2K(\cos\theta)}$ and conjectured that the maximum value 
of $(1-x^2)^2 \abs{S_{f_\alpha}(x)}$ occurs
at $x=0$.
(Note the typographical error on the second $\alpha$ in the original 
statement in~\cite{PM00}.)


We claim that maximizing $(1-x^2)^2 \abs{S_f(x)}$ 
is equivalent to maximizing $(1-x^2)^2 {S_f(x)}$, i.e., we claim that
$\max (1-x^2)^2 {S_f(x)} > -\min (1-x^2)^2 {S_f(x)}$.  We will show that

\begin{equation} \label{claim_min}
\min (1-x^2)^2 {S_f(x)} \ge -2
\end{equation}

\vspace{0.1in}

\noindent over $ 0 < x < 1$ and $0<\theta<\pi/2$. 
Then, as Mej\'ia and Pommerenke~\cite{PM00} noted (and as we will show)
that $\max (1-x^2)^2 {S_f(x)} > 2$, the claim will hold.  

We simplify our algebraic computations
by replacing $x$ with $\sqrt{x}$ in the above expression and obtain

\begin{equation} \label{E:m}
(1-x)^{2} \frac{2( (\cos(2\theta)+\alpha^2) 
+ (\cos^2(2\theta)-2\alpha^2\cos(2 \theta) -3)x 
+ (\cos(2\theta) +\alpha^2)x^2)}{(1-2\cos(2\theta)x +x^2)^2}
\end{equation}

\noindent where $ 0 \le x < 1$ and $0<\theta<\pi/2$.

To minimize (\ref{E:m}), we note that the coefficient of $\alpha^2$ is $1-2\cos(2\theta)x +x^2$ which
is nonnegative.  Hence, substituting $\alpha=0$ into (\ref{E:m}), we have that (\ref{E:m}) is bounded below by

\begin{equation} \label{E:m1}
(1-x)^{2} \frac{2( \cos(2\theta) 
+ (\cos^2(2\theta) -3)x 
+ \cos(2\theta) x^2)}{(1-2\cos(2\theta)x +x^2)^2}
\end{equation}

We will show that (\ref{E:m1}) is bounded below by $-2$ by showing that 

\begin{align}\notag
&(1-x)^{2} \frac{2( \cos(2\theta) 
+ (\cos^2(2\theta) -3)x 
+ \cos(2\theta) x^2)}{(1-2\cos(2\theta)x +x^2)^2} + 2 \\ \notag
&= \frac{p(c,x)}{(1-2\cos(2\theta)x +x^2)^2} > 0
\end{align}

\noindent where

$$
 p(c,x) = (2x+4x^2 + 2x^3)c^2 + (2-12x+4x^2-12x^3+2x^4)c + (2-6x+16x^2-6x^3+2x^4)
$$

\noindent and $c = \cos(2\theta)$ satistifies $-1 < c < 1$.  

It is sufficient to show that $p(c,x) > 0$ for $-1 < c < 1$ and $0 < x < 1$.
We note that $p(c,x)$ is quadratic in $c$ for each fixed $x$ and the coefficent 
of $c^2$ is positive.  Let 

$$
c(x) = -\frac{2-12x+4x^2-12x^3+2x^4}{2(2x+4x^2 + 2x^3)},
$$

\noindent then $c(x)$ is the location of the vertex of the quadratic $p(c,x)$ for each fixed $x$.
We claim that $c(x)$ is an increasing function of $x$ on $(0,1)$.  To see this, we note that

$$c'(x) = \frac{n(x)}{2(x+1)(1+2x+x^2)x^2}$$

\noindent where $n(x) = -x^5-3x^4+14x^3-14x^2+3x+1$.  Clearly, $n(0) = 1$ and $n(1) = 0$.  A Sturm sequence argument shows that $n(x)$ has exactly one zero on $(0,1]$.  Hence, we have $c'(x) > 0$ on 
$(0,1)$ and $c(x)$ is increasing on $(0,1)$.

Let $x_0$ be the unique solution of $c(x) = -1$ 
for $0 < x < 1$, $x_0 \approx 0.1197$. 
For $0 < x < x_0$, the location $c(x)$ of the vertex of $p(c,x)$ satisfies $c(x) < -1$.  
Hence over $-1 < c < 1$ and $0 < x < 1$, 
we have $ \min p(c,x) > p(-1,x) = 8x + 16x^2 +8x^3 > 0$.  For $x_0 < x < 1$, we have
$-1 < c(x) < 1$ and, hence, 

$$ \min p(c,x) = p(c,x)|_{c=c(x)} = \frac{q(x)}{2x(1+2x+x^2)} $$

\noindent where

$$ 
q(x) = -x^8 + 16x^7 - 44x^6 + 48x^5 - 3xx^4 +48x^3-44x^2+16x -1.
$$

Clearly, $q(1) = 0$ and $q(1/2)=\frac{143}{256} > 0$.  A Sturm sequence argument shows that on 
the interval $(x_0,1]$ that $q$ has exactly one zero.  Hence, on $(0,1)$, $q(x) > 0$.

Thus, $p(c,x) > 0$ for 
$-1 < c < 1$ and $0 < x < 1$ and, consequently, $\min (1-x^2)^2 {S_f(x)} \ge -2$.

To maximize (\ref{E:m}), we substitute $\alpha=\frac{\pi}{2 K(\cos(\theta))}$,
take the derivative with respect to $x$, and simplify to obtain

\begin{align}
\frac{(x^2-1)(\cos(2\theta) -1)}
{K^2(\cos(\theta))(x^2-2\cos(2\theta)x+1)^3}
&(K^2(\cos\theta)(-10x^2 \cos(2\theta) \\\notag
&-6x^2 
+12x\cos(2\theta)-4x\cos^2(2\theta)\\\notag
&+24x-10\cos(2\theta)-6) -\pi^2(x^2-2\cos(2\theta)+1))
\end{align}


Notice that the denominator of this expression is always positive.
%since 
%\[
%x^2-2\cos(2\theta)x+1 > x^2-2x+1 =(x-1)^2 >0.
%\]
On the other hand, the numerator changes sign as $x$ varies.  However,
we will see that for fixed $\theta$, the numerator can change sign
only once.

To justify this claim, fix $\theta$ and let 
\begin{align} \notag
A(x)=K^2(\cos\theta)(&-10x^2 \cos(2\theta) -6x^2 
+12x\cos(2\theta)-4x\cos^2(2\theta)\\\notag
& +24x-10\cos(2\theta)-6)
-\pi^2(x^2-2\cos(2\theta)+1).
\end{align}
Notice that $(x^2-1)(\cos(2\theta)-1)>0$ for $0\leq x<1$ and
$0<\theta<\pi/2$, so we need
only show that $A(x)$ changes sign only once.

It follows from known properties~\cite{AVV97} of $K$ that 
$K(\cos(\theta))>\pi/2$; thus, since the coefficient of $K(\cos(\theta))$
in \eqref{E:ktheta} is positive,
\begin{align}\label{E:ktheta}
A'(x) & = K^2(\cos(2\theta))(-20x\cos(2\theta) -12x+12\cos(2\theta) - 4\cos^2(2\theta) + 24) - 2\pi^2(x-\cos(2\theta))\\\notag
&>\left(\frac{\pi}{2}\right)^2(-20x\cos(2\theta) 
-12x+12\cos(2\theta) - 4\cos^2(2\theta) + 24) - 2\pi^2(x-\cos(2\theta))\\\notag
&= (-10\pi^2\cos^2\theta)x + (14\pi^2\cos^2\theta-4\pi^2\cos^4\theta )
\end{align}


The right hand side of the above inequality is linear in $x$ and positive
at both $x=0$ and $x=1$.  Thus $A'(x)>0$  and
$A(x)$ is increasing for all $0\leq x <1$.  Consequently, $A(x)$
can change sign at most once.

Notice
\[
A(0)=-10\cos(2\theta)K^2(\cos(2\theta)) -6K^2(\cos(\theta))-\pi^2 <0
\]
for all $0<\theta<\pi/2$.  

Thus, as a consequence of the above argument, for each fixed $\theta$,
$ %\frac\partial}{\partial x} 
(1-x^2)^2S_{f_\alpha}(x)$ is decreasing at $x=0$ and
switches from decreasing to increasing at most once on the interval
$0\leq x <1$. (Recall that to establish this fact, we replaced
$x$ by $\sqrt{x}$ in our calculations.  We now  return $x$ to its
original meaning.)

Clearly  $(1-x^2)^2S_{f_\alpha}(x)$ equals $0$ when $x=1$.  
Thus even after $(1-x^2)^2S_{f_\alpha}(x)$ begins increasing, it never
becomes larger than $0$.

On the other hand, when $x=0$,
\begin{equation}\notag
(1-x^2)^2S_{f_\alpha}(x)|_{x=0} = 2\cos(2\theta)+\frac{\pi^2}{2K^2(\cos(\theta))}.
\end{equation}
Applying Lemma~\ref{L:matti} with $c=1/2$ %from~[1,~p.~54],
we obtain
\begin{align}\notag
(1-x^2)^2S_{f_\alpha}(x)|_{x=0} 
&\geq  2\cos(2\theta)+\frac{\pi^2}{\frac{\pi^2}{2
\sqrt{1-\cos^2(2\theta)}}} \\\notag
&=2(1+2\sin(\theta))(1-\sin(\theta))>0
\end{align}
for all $0<\theta<\pi/2$.
Thus, $\sup_{0\leq x<1}(1-x^2)^2 S_{f_\alpha}(x)$ is indeed achieved for $x=0$.
See Figure~\ref{F:x-running}.



\begin{figure}
\begin{center}\scalebox{0.45}{
%\begin{rotate}{-90}
%\begin{minipage}{\linewidth}
\epsfig{file=x-running,angle=-90} }
%\end{minipage}
%\end{rotate} }
\caption{The graph of $(1-x^2)^2{S_{f_\alpha}(x)}$ for $0\leq x<1$ and
$\alpha=0.5343598$ ($\theta=0.218202$).}
\label{F:x-running}
\end{center}
\end{figure}

Now the only question is to determine for which value (or values) of $\theta$ 
\[
\norm{S_{f_\alpha}}_\D=S_{f_\alpha}(0) 
= 4\cos^2(\theta)-2+\frac{\pi^2}{2K^2(\cos(\theta))}
\]
is maximized.  
To this end, let $s=\cos \theta$ and let 
\[
g(s)=4s^2-2+\frac{\pi^2}{2K^2({s})}.
\]

Notice that 
\[
\frac{dg}{d\theta}=\frac{dg}{ds}\frac{ds}{d\theta},
\]
and $\frac{ds}{d\theta}=-\sin(\theta)<0$ for $0<\theta<\pi/2$.
Thus since $\cos(\theta)$ is decreasing for $0<\theta<\pi/2$, 
if $\frac{dg}{ds}$ switches signs at most once, then $\norm{S_{f_\alpha}}_\D$
will have a relative maximum  for  at most one value of $\theta$.

Computing this derivative and simplifying, we obtain
\begin{equation}\label{E:dg}
\frac{dg}{ds}(s)= 8s - 
\frac{\pi^2}{s}\frac{
E\left(s\right)-(1-s^2)K( s))}{(1-s^2)K^3({s})}.
\end{equation}

For $0<s<1$, we rewrite\eqref{E:dg} as
\[
%4-\frac{\pi^2}{2}\frac{
%\frac{E(t)-(1-t^2)K(t)}{t^2}}
%{(1-t^2)K^3(t)} =
2s\left( 4-\frac{\pi^2}{2}\frac{E(s)-(1-s^2)K(s)}{s^2}
\frac{1}{\left((\sqrt{1-s^2})^{2/3} K(s)\right)^{3}}\right).
\]



Applying the the first part of Lemma~\ref{L:matti} with $c=2/3$,
we observe that
\[
\frac{1}{\left((\sqrt{1-s^2})^{2/3} K(s)\right)^{3}}
\]
is increasing for $0<s<1$.
Similarly, by the second part of Lemma~\ref{L:matti}, 
\[
\frac{E(s)-(1-s^2)K(s)}{s^2}
\]
is increasing for $0<s<1$.

Consequently, $\frac{dg}{ds}$ can be 
written as  a product of $2s$, which is positive, and a factor, which is 
$4$ minus an increasing function.  Thus, $\frac{dg}{ds}$ has at most one zero.
Moreover, Lemma~\ref{L:matti} also implies that
\begin{align}
\lim_{s\to 0^+}\frac{E(s)-(1-s^2)K(s)}{s^2} &=\frac{\pi}{4} \\\notag 
\lim_{s\to 0^+} {\left((\sqrt{1-s^2})^{2/3} K(s)\right)^{3}}
&= \left(\frac{\pi}{2}\right)^{3}.
\end{align}
Thus,
\[
\lim_{s\to 0^+} \frac{dg}{ds} = 4-\frac{\pi^2}{2}\frac{\pi/4}{(\pi/2)^3} =3.
\]

On the other hand, 
\begin{align}
\lim_{s\to 1^-}\frac{E(s)-(1-s^2)K(s)}{s^2} &= 1 \\\notag
\lim_{s\to 1^-} {\left((\sqrt{1-s^2})^{2/3} K(s)\right)^{3}}
&= 0.
\end{align}
Hence,
\[
\lim_{s\to 1^-}\frac{dg}{ds} = -\infty.
\]



Since $\frac{dg}{ds}$ 
decreases from $3$ to minus infinity as $s$ increases from $0$ to $1$,
 we have proven that
$\norm{S_{f_\alpha}}_\D$ does have a unique maxiumum.  
A Maple computation shows
that this maximum is approximately $2.383635$ 
and occurs for $\theta\approx 0.218202$,
which corresponds to $\alpha\approx 0.5343598$.  See Figure~\ref{F:t-running}.



\begin{figure}
\begin{center}\scalebox{0.45}{
\epsfig{file=t-running} }
\caption{The graph of $\norm{S_{f_\alpha}}_\D$ for  $0<\theta<\pi/2$.}
\label{F:t-running}
\end{center}
\end{figure}

\end{proof}

\begin{remark}
Mej\'ia and Pommerenke~\cite{MPV01}  have recently determined the 
sharp upper bound 
for $\abs{f'(re^{i\theta})}$, $0\leq r <1$, for
 hyperbolically convex functions $f$.
Applying the techniques of this paper we give another proof and
extensively generalize this result~\cite{BOP03}.
We have also answered in~\cite{BPW04} several other related conjectures of 
Pommerenke~\cite{pomm1}.


\end{remark}

%\bibliographystyle{amsplain} \bibliography{hc,roger,eqcp,dis}

\def\cprime{$'$}
\providecommand{\bysame}{\leavevmode\hbox to3em{\hrulefill}\thinspace}
\providecommand{\MR}{\relax\ifhmode\unskip\space\fi MR }
% \MRhref is called by the amsart/book/proc definition of \MR.
\providecommand{\MRhref}[2]{%
  \href{http://www.ams.org/mathscinet-getitem?mr=#1}{#2}
}
\providecommand{\href}[2]{#2}


\begin{thebibliography}{10000}

\bibitem{AVV97}
G.D. Anderson, M.K. Vamanamurthy, and M.K. Vuorinen, \emph{Conformal
  invariants, inequalities, and quasiconformal mappings}, J. Wiley and Sons,
  1997.

\bibitem{BPW04}
R.~W. Barnard, K.~Pearce, and G.~Brock Williams, \emph{Three extremal problems
  for hyperbolically convex functions}, Computational Methods and Function
  Theory \textbf{4} (2004), no.~1, 97--109.

\bibitem{BL75}
Roger Barnard and John~L. Lewis, \emph{Subordination theorems for some classes
  of starlike fumctions}, Pacific J. Math. \textbf{56} (1975), no.~2, 333--366.
  \MR{52 \#721}

\bibitem{BOP03}
Roger~W. Barnard, G.~Lenny Ornas, and Kent Pearce, \emph{A variational method
  for hyperbolically convex functions}, Complex Variables and Elliptic Equations 
  \textbf{51} (2006), no.~4, 313-327.

\bibitem{aB83}
A.~Beardon, \emph{Geometry of discrete groups}, Springer-Verlag, Berlin, 1983.

\bibitem{GL2000}
Frederick~P. Gardiner and Nikola Lakic, \emph{Quasiconformal {T}eichm\"uller
  theory}, Mathematical Surveys and Monographs, vol.~76, American Mathematical
  Society, 2000.

\bibitem{oL87}
O.~Lehto, \emph{Univalent functions and {T}eichm\"uller spaces},
  Springer-Verlag, Berlin - Heidelberg - New York, 1987.

\bibitem{MM94}
William Ma and David Minda, \emph{Hyperbolically convex functions}, Ann. Polon.
  Math. \textbf{60} (1994), no.~1, 81--100. \MR{95k:30037}

\bibitem{MM99}
\bysame, \emph{Hyperbolically convex functions. {II}}, Ann. Polon. Math.
  \textbf{71} (1999), no.~3, 273--285. \MR{2000j:30020}

\bibitem{PM00}
Diego Mej\'ia and Christian Pommerenke, \emph{On hyperbolically convex
  functions}, J. Geom. Anal. \textbf{10} (2000), no.~2, 365--378.

\bibitem{MP00b}
\bysame, \emph{On spherically convex univalent functions}, Michigan Math. J.
  \textbf{47} (2000), no.~1, 163--172. \MR{2001a:30013}

\bibitem{MP01}
\bysame, \emph{Sobre la derivada {Schawarziana} de aplicaciones conformes
  hiperb\'olicamente}, Revista Colombiana de Matem\'aticas \textbf{35} (2001),
  no.~2, 51--60.

\bibitem{MP02}
\bysame, \emph{Hyperbolically convex functions, dimension and capacity},
  Complex Var. Theory Appl. \textbf{47} (2002), no.~9, 803--814.
  \MR{2003f:30017}

\bibitem{MP00}
\bysame, \emph{On the derivative of hyperbolically convex functions}, Ann.
  Acad. Sci. Fenn. Math. \textbf{27} (2002), no.~1, 47--56.

\bibitem{MPV01}
Diego Mej\'ia, Christian Pommerenke, and Alexander Vasil{\cprime}ev,
  \emph{Distortion theorems for hyperbolically convex functions}, Complex
  Variables Theory Appl. \textbf{44} (2001), no.~2, 117--130. \MR{2003e:30025}

\bibitem{zN49}
Zeev Nehari, \emph{The {S}chwarzian derivative and schlicht functions}, Bull.
  Amer. Math. Soc. \textbf{55} (1949), 545--551. \MR{10,696e}

\bibitem{zN52}
\bysame, \emph{Conformal mapping}, Dover, 1952.

\bibitem{zN76}
\bysame, \emph{A property of convex conformal maps}, J. Analyse Math.
  \textbf{30} (1976), 390--393. \MR{55 \#12901}

\bibitem{pomm1}
Christian Pommerenke, private communication.

\bibitem{cP92}
Ch. Pommerneke, \emph{Boundary behavior of conformal maps}, Springer - Verlag,
  Berlin - Heidelberg - New York, 1992.

\bibitem{aS98}
A.~Yu. Solynin, \emph{Some extremal problems on the hyperbolic polygons},
  Complex Variables Theory Appl. \textbf{36} (1998), no.~3, 207--231.
  \MR{99j:30030}

\bibitem{aS99}
\bysame, \emph{Moduli and extremal metric problems}, Algebra i Analiz
  \textbf{11} (1999), no.~1, 3--86, translation in St. Petersburg Math. J.
  \textbf{11} (2000), no. 1, 1--65. \MR{2001b:30058}

\end{thebibliography}



\affiliationone{% in this example, two authors share an institution
   Roger W. Barnard, Kent Pearce and G. Brock Williams\\
   Department of Mathematics and Statistics, Texas Tech University,
 Lubbock, Texas  79409\\United States
   \email{roger.w.barnard@ttu.edu\\kent.pearce@ttu.edu\\brock.williams@ttu.edu}}
\affiliationtwo{% in this example, one author has two addresses}
   Leah Cole\\
   Department of Mathematics, Wayland Baptist University, 
 Plainview, Texas  79072\\United States
   \email{colel@wbu.edu}}


\end{document}


